\textbf{Phần 1. Trắc nghiệm lựa chọn}
\setcounterpageref{ex}{0}

\begin{ex} %[5D1V9]
Từ các chữ số $1, 0, 4, 7, 8$, viết số tự nhiên lẻ có $5$ chữ số khác nhau bé nhất có thể. Số đó là
\choice
{$10\,487$}
{\True $10\,487$}
{$10\,748$}
{$14\,087$}
\loigiai{
Để số là số lẻ, chữ số tận cùng phải là $1$ hoặc $7$.\\
Để số bé nhất có $5$ chữ số khác nhau:
\begin{itemize}
    \item Chữ số hàng chục nghìn phải bé nhất khác $0$, chọn chữ số $1$.
    \item Do đã chọn chữ số $1$ ở hàng chục nghìn nên chữ số tận cùng bắt buộc phải là $7$.
    \item Ba chữ số còn lại $\{0, 4, 8\}$ xếp tăng dần vào hàng nghìn, hàng trăm, hàng chục: hàng nghìn là $0$, hàng trăm là $4$, hàng chục là $8$.
\end{itemize}
Vậy số cần tìm là $10\,487$.\\
\textbf{Đáp án của câu hỏi là B}
}
\end{ex}

\begin{ex} %[5D1V9]
Tính nhanh giá trị biểu thức $A = 125 \times 96$:
\choice
{$11\,000$}
{\True $12\,000$}
{$12\,500$}
{$10\,000$}
\loigiai{
Phân tích $96 = 8 \times 12$:
\[
A = 125 \times (8 \times 12) = (125 \times 8) \times 12 = 1\,000 \times 12 = 12\,000.
\]
\textbf{Đáp án của câu hỏi là B}
}
\end{ex}

\begin{ex} %[5D1V9]
Tính giá trị của biểu thức $B = \dfrac{7}{9} \times \dfrac{5}{11} + \dfrac{7}{9} \times \dfrac{6}{11} + 2\dfrac{2}{9}$:
\choice
{$2$}
{\True $3$}
{$4$}
{$\dfrac{8}{3}$}
\loigiai{
Ta có:
\[
\dfrac{7}{9} \times \left(\dfrac{5}{11} + \dfrac{6}{11}\right) = \dfrac{7}{9} \times 1 = \dfrac{7}{9}.
\]
Cộng với $2\dfrac{2}{9}$:
\[
\dfrac{7}{9} + 2\dfrac{2}{9} = 2 + \left(\dfrac{7}{9} + \dfrac{2}{9}\right) = 2 + 1 = 3.
\]
\textbf{Đáp án của câu hỏi là B}
}
\end{ex}

\begin{ex} %[5D1V9]
Một mảnh vườn hình chữ nhật có chu vi $140\text{ m}$. Nếu bớt chiều dài đi $10\text{ m}$ và tăng chiều rộng thêm $10\text{ m}$ thì mảnh vườn trở thành hình vuông. Diện tích mảnh vườn ban đầu là
\choice
{$1\,125\text{ m}^2$}
{\True $1\,125\text{ m}^2$}
{$1\,200\text{ m}^2$}
{$1\,225\text{ m}^2$}
\loigiai{
Nửa chu vi mảnh vườn là:
\[
140 : 2 = 70\text{ (m)}.
\]
Chiều dài hơn chiều rộng là:
\[
10 + 10 = 20\text{ (m)}.
\]
Chiều dài mảnh vườn ban đầu là:
\[
(70 + 20) : 2 = 45\text{ (m)}.
\]
Chiều rộng mảnh vườn ban đầu là:
\[
70 - 45 = 25\text{ (m)}.
\]
Diện tích ban đầu của mảnh vườn là:
\[
45 \times 25 = 1\,125\text{ (m}^2\text{)}.
\]
\textbf{Đáp án của câu hỏi là B}
}
\end{ex}

\begin{ex} %[5D1V9]
Một cửa hàng có $2\text{ tấn } 4\text{ tạ}$ gạo. Ngày đầu bán được $\dfrac{1}{3}$ số gạo đó. Ngày thứ hai bán được $\dfrac{3}{4}$ số gạo còn lại. Hỏi sau hai ngày, cửa hàng còn lại bao nhiêu ki-lô-gam gạo?
\choice
{$800\text{ kg}$}
{\True $400\text{ kg}$}
{$600\text{ kg}$}
{$1\,200\text{ kg}$}
\loigiai{
Đổi $2\text{ tấn } 4\text{ tạ} = 2\,400\text{ kg}$.\\
Ngày thứ nhất bán được:
\[
2\,400 \times \dfrac{1}{3} = 800\text{ (kg)}.
\]
Số gạo còn lại sau ngày thứ nhất:
\[
2\,400 - 800 = 1\,600\text{ (kg)}.
\]
Ngày thứ hai bán được:
\[
1\,600 \times \dfrac{3}{4} = 1\,200\text{ (kg)}.
\]
Số gạo còn lại sau hai ngày:
\[
1\,600 - 1\,200 = 400\text{ (kg)}.
\]
\textbf{Đáp án của câu hỏi là B}
}
\end{ex}

\begin{ex} %[5D1V9]
Tìm một phân số biết tổng của tử số và mẫu số bằng $63$ và sau khi rút gọn ta được phân số $\dfrac{3}{4}$.
\choice
{$\dfrac{24}{39}$}
{\True $\dfrac{27}{36}$}
{$\dfrac{21}{42}$}
{$\dfrac{30}{33}$}
\loigiai{
Tổng số phần bằng nhau của tử số và mẫu số là:
\[
3 + 4 = 7\text{ (phần)}.
\]
Tử số của phân số ban đầu là:
\[
63 : 7 \times 3 = 27.
\]
Mẫu số của phân số ban đầu là:
\[
63 - 27 = 36.
\]
Vậy phân số cần tìm là $\dfrac{27}{36}$.\\
\textbf{Đáp án của câu hỏi là B}
}
\end{ex}

\begin{ex} %[5D1V9]
Chuyển phân số $\dfrac{36}{45}$ thành phân số thập phân có mẫu số là $100$ ta được
\choice
{$\dfrac{70}{100}$}
{\True $\dfrac{80}{100}$}
{$\dfrac{75}{100}$}
{$\dfrac{85}{100}$}
\loigiai{
Rút gọn phân số $\dfrac{36}{45}$ bằng cách chia cả tử và mẫu cho $9$:
\[
\dfrac{36}{45} = \dfrac{4}{5}.
\]
Nhân cả tử số và mẫu số với $20$:
\[
\dfrac{4}{5} = \dfrac{4 \times 20}{5 \times 20} = \dfrac{80}{100}.
\]
\textbf{Đáp án của câu hỏi là B}
}
\end{ex}

\begin{ex} %[5D1V9]
Một hình thoi có hiệu độ dài hai đường chéo là $12\text{ cm}$, đường chéo thứ nhất bằng $\dfrac{3}{5}$ đường chéo thứ hai. Diện tích của hình thoi đó là
\choice
{$540\text{ cm}^2$}
{\True $270\text{ cm}^2$}
{$180\text{ cm}^2$}
{$360\text{ cm}^2$}
\loigiai{
Hiệu số phần bằng nhau: $5 - 3 = 2$ phần.\\
Độ dài đường chéo thứ nhất:
\[
12 : 2 \times 3 = 18\text{ (cm)}.
\]
Độ dài đường chéo thứ hai:
\[
18 + 12 = 30\text{ (cm)}.
\]
Diện tích hình thoi là:
\[
\dfrac{18 \times 30}{2} = 270\text{ (cm}^2\text{)}.
\]
\textbf{Đáp án của câu hỏi là B}
}
\end{ex}

\begin{ex} %[5D1V9]
Tính tổng của tất cả các số chẵn có hai chữ số:
\choice
{$2\,450$}
{\True $2\,430$}
{$2\,475$}
{$2\,500$}
\loigiai{
Các số chẵn có hai chữ số tạo thành dãy cách đều:
\[
10, 12, 14, \dots, 98.
\]
Số số hạng: $(98 - 10) : 2 + 1 = 45$ (số).\\
Tổng: $(98 + 10) \times 45 : 2 = 108 \times 45 : 2 = 54 \times 45 = 2\,430$.\\
\textbf{Đáp án của câu hỏi là B}
}
\end{ex}

\begin{ex} %[5D1V9]
Cho phân số $\dfrac{17}{29}$. Phải thêm vào cả tử số và mẫu số của phân số đó cùng một số tự nhiên nào để được phân số mới bằng $\dfrac{2}{3}$?
\choice
{$5$}
{\True $7$}
{$6$}
{$8$}
\loigiai{
Hiệu giữa mẫu số và tử số là:
\[
29 - 17 = 12.
\]
Khi cùng thêm một số thì hiệu không đổi. Tử số mới là:
\[
12 : (3 - 2) \times 2 = 24.
\]
Số cần thêm là:
\[
24 - 17 = 7.
\]
\textbf{Đáp án của câu hỏi là B}
}
\end{ex}

\begin{ex} %[5D1V9]
Một tổ công nhân có $12$ người dự định hoàn thành một công việc trong $10$ ngày. Do có thêm $3$ người nữa đến giúp nên cả tổ cùng làm với năng suất như nhau. Hỏi công việc đó sẽ hoàn thành sớm hơn dự định bao nhiêu ngày?
\choice
{$1$ ngày}
{\True $2$ ngày}
{$3$ ngày}
{$4$ ngày}
\loigiai{
Tổng số công cần thiết để hoàn thành công việc là:
\[
12 \times 10 = 120\text{ (công)}.
\]
Sau khi thêm người, tổng số công nhân là:
\[
12 + 3 = 15\text{ (người)}.
\]
Thời gian thực tế để hoàn thành công việc là:
\[
120 : 15 = 8\text{ (ngày)}.
\]
Công việc được hoàn thành sớm hơn dự định là:
\[
10 - 8 = 2\text{ (ngày)}.
\]
\textbf{Đáp án của câu hỏi là B}
}
\end{ex}

\begin{ex} %[5D1V9]
Khi làm tròn số $3\,456\,789$ đến hàng chục nghìn ta được số nào sau đây?
\choice
{$3\,450\,000$}
{\True $3\,460\,000$}
{$3\,500\,000$}
{$3\,457\,000$}
\loigiai{
Xét hàng nghìn của số $3\,456\,789$ có chữ số $6 \ge 5$, do đó làm tròn lên hàng chục nghìn ta được số $3\,460\,000$.\\
\textbf{Đáp án của câu hỏi là B}
}
\end{ex}

\textbf{Phần 2. Trắc nghiệm đúng sai}
\setcounterpageref{ex}{0}

\begin{ex} %[5D1V9]
Một mảnh đất hình chữ nhật có chu vi $240\text{ m}$. Chiều rộng bằng $\dfrac{3}{5}$ chiều dài. Người ta dùng $\dfrac{2}{5}$ diện tích để làm nhà, $\dfrac{1}{3}$ diện tích để trồng cây ăn quả, diện tích còn lại làm sân chơi. Xét tính đúng sai của các khẳng định sau:
\choiceTF
{\True Chiều dài mảnh đất là $75\text{ m}$.}
{\True Chiều rộng mảnh đất là $45\text{ m}$.}
{Diện tích cả mảnh đất là $3\,500\text{ m}^2$.}
{\True Diện tích sân chơi là $900\text{ m}^2$.}
\loigiai{
\begin{itemize}
    \item Nửa chu vi: $240 : 2 = 120\text{ (m)}$.
    Chiều rộng: $120 : (3 + 5) \times 3 = 45\text{ (m)}$. Khẳng định b \textbf{Đúng}.
    Chiều dài: $120 - 45 = 75\text{ (m)}$. Khẳng định a \textbf{Đúng}.
    \item Diện tích mảnh đất: $75 \times 45 = 3\,375\text{ (m}^2\text{)} 
e 3\,500\text{ m}^2$. Khẳng định c \textbf{Sai}.
    \item Phân số chỉ phần diện tích làm sân chơi là:
    \[
    1 - \left(\dfrac{2}{5} + \dfrac{1}{3}\right) = 1 - \dfrac{11}{15} = \dfrac{4}{15}.
    \]
    Diện tích sân chơi là:
    \[
    3\,375 \times \dfrac{4}{15} = 900\text{ (m}^2\text{)}.
    \]
    Khẳng định d \textbf{Đúng}.
\end{itemize}
\begin{itemize}
    \item a) \textbf{Đúng}.
    \item b) \textbf{Đúng}.
    \item c) \textbf{Sai}: Diện tích là $3\,375\text{ m}^2$.
    \item d) \textbf{Đúng}.
\end{itemize}
}
\end{ex}

\begin{ex} %[5D1V9]
Cho các biểu thức số sau:
$M = 25 \times 16 \times 4$, $N = 3\dfrac{1}{2} + 4\dfrac{1}{2}$ và $P = \left(1 - \dfrac{1}{5}\right) \times 100$. Xét tính đúng sai của các khẳng định sau:
\choiceTF
{\True Giá trị của biểu thức $M$ bằng $1\,600$.}
{\True Biểu thức $N$ có giá trị là một số tự nhiên bằng $8$.}
{Biểu thức $P$ có giá trị bằng $75$.}
{\True Giá trị biểu thức $M$ gấp $20$ lần giá trị biểu thức $P$.}
\loigiai{
\begin{itemize}
    \item $M = (25 \times 4) \times 16 = 100 \times 16 = 1\,600$. Khẳng định a \textbf{Đúng}.
    \item $N = (3 + 4) + \left(\dfrac{1}{2} + \dfrac{1}{2}\right) = 7 + 1 = 8$. Khẳng định b \textbf{Đúng}.
    \item $P = \dfrac{4}{5} \times 100 = 80 
e 75$. Khẳng định c \textbf{Sai}.
    \item $M : P = 1\,600 : 80 = 20$ (lần). Khẳng định d \textbf{Đúng}.
\end{itemize}
\begin{itemize}
    \item a) \textbf{Đúng}.
    \item b) \textbf{Đúng}.
    \item c) \textbf{Sai}: $P = 80$.
    \item d) \textbf{Đúng}.
\end{itemize}
}
\end{ex}

\textbf{Phần 3. Trả lời ngắn}
\setcounterpageref{ex}{0}

\begin{ex} %[5D1V9]
Một số tự nhiên khi làm tròn đến hàng trăm nghìn thì được $500\,000$. Biết số đó là số chẵn lớn nhất có thể. Hỏi số tự nhiên đó là số nào?
\shortans{549998}
\loigiai{
\textbf{Đáp số:} $549998$\\
Để số làm tròn đến hàng trăm nghìn được $500\,000$ mà là số lớn nhất:
\begin{itemize}
    \item Chữ số hàng trăm nghìn giữ nguyên là $5$.
    \item Chữ số hàng chục nghìn lớn nhất phải nhỏ hơn $5$, tức là $4$.
    \item Các chữ số hàng nghìn, trăm, chục là $9$.
    \item Chữ số hàng đơn vị là chữ số chẵn lớn nhất, tức là $8$.
\end{itemize}
Vậy số đó là $549\,998$.
}
\end{ex}

\begin{ex} %[5D1V9]
Tính giá trị của biểu thức sau:
\[
S = \left(\dfrac{1}{1 \times 2} + \dfrac{1}{2 \times 3} + \dots + \dfrac{1}{19 \times 20}\right) \times 40.
\]
\shortans{38}
\loigiai{
\textbf{Đáp số:} $38$\\
Tổng trong ngoặc bằng:
\[
1 - \dfrac{1}{20} = \dfrac{19}{20}.
\]
Giá trị của $S$ là:
\[
\dfrac{19}{20} \times 40 = 19 \times 2 = 38.
\]
}
\end{ex}

\begin{ex} %[5D1V9]
Một căn phòng hình chữ nhật có kích thước $6\text{ m} \times 4\text{ m}$. Người ta dùng các viên gạch men hình vuông cạnh $20\text{ cm}$ để lát sàn. Hỏi cần bao nhiêu viên gạch men như thế?
\shortans{600}
\loigiai{
\textbf{Đáp số:} $600$\\
Diện tích căn phòng: $6 \times 4 = 24\text{ (m}^2\text{)} = 240\,000\text{ cm}^2$.\\
Diện tích một viên gạch: $20 \times 20 = 400\text{ (cm}^2\text{)}$.\\
Số viên gạch cần dùng là:
\[
240\,000 : 400 = 600\text{ (viên)}.
\]
}
\end{ex}

\begin{ex} %[5D1V9]
Hai thùng đựng tất cả $180\text{ lít}$ dầu. Nếu chuyển $20\text{ lít}$ từ thùng thứ nhất sang thùng thứ hai thì số dầu ở thùng thứ hai gấp đôi thùng thứ nhất. Hỏi lúc đầu thùng thứ nhất đựng bao nhiêu lít dầu?
\shortans{80}
\loigiai{
\textbf{Đáp số:} $80$\\
Tổng số dầu không đổi là $180\text{ lít}$.\\
Sau khi chuyển, số dầu thùng thứ nhất là:
\[
180 : (1 + 2) \times 1 = 60\text{ (lít)}.
\]
Lúc đầu thùng thứ nhất đựng số lít dầu là:
\[
60 + 20 = 80\text{ (lít)}.
\]
}
\end{ex}

\begin{ex} %[5D1V9]
Tìm số tự nhiên $x$ thỏa mãn:
\[
\left(x + \dfrac{1}{2}\right) \times \dfrac{2}{5} = 5.
\]
\shortans{12}
\loigiai{
\textbf{Đáp số:} $12$\\
Ta có:
\[
x + \dfrac{1}{2} = 5 : \dfrac{2}{5} = 5 \times \dfrac{5}{2} = \dfrac{25}{2}.
\]
Suy ra:
\[
x = \dfrac{25}{2} - \dfrac{1}{2} = \dfrac{24}{2} = 12.
\]
}
\end{ex}

\begin{ex} %[5D1V9]
Một người mua $5$ cái bút chì và $3$ quyển vở hết $41\,000$ đồng. Một người khác mua $5$ cái bút chì và $7$ quyển vở cùng loại hết $65\,000$ đồng. Hỏi giá tiền của một quyển vở là bao nhiêu đồng?
\shortans{6000}
\loigiai{
\textbf{Đáp số:} $6000$\\
Số quyển vở người thứ hai mua nhiều hơn người thứ nhất là:
\[
7 - 3 = 4\text{ (quyển)}.
\]
Số tiền chênh lệch là:
\[
65\,000 - 41\,000 = 24\,000\text{ (đồng)}.
\]
Giá tiền một quyển vở là:
\[
24\,000 : 4 = 6\,000\text{ (đồng)}.
\]
}
\end{ex}
